\begin{abstract}
Systems that repair language-model-generated code by feeding verifier output back into the prompt must choose what goes in the verify stage, and that choice currently rests on an untested intuition: that a more formally rigorous verifier produces more precise feedback and therefore better repair. Prior work studies one feedback category at a time, on different backbones, and reports what repair achieved but never what it cost. We hold the model, the benchmark, the prompt template, and the repair budget fixed while varying only the verifier across four categories --- execution monitoring, static analysis, type checking, and SMT solving --- over 421 HumanEval and MBPP problems, and we measure wall-clock overhead alongside correctness. The intuition inverts. Repair success falls monotonically as feedback grows longer, giving a rank correlation of exactly $-1.0$: verbose static-analysis diagnostics fix fewer programs than short execution tracebacks. Cost inverts the ranking again, because static analysis runs an estimated seventeen times faster and returns an estimated sixteen times more correctness per second (overhead from calibrated mock run; ordering robust, magnitudes are estimates --- see \S4.4), while execution monitoring delivers twice the absolute improvement. Repair appears bounded by the model's ability to extract one actionable signal, not by how much the feedback says.
\end{abstract}
